\scanpage{46}
\margin{45}
- \mspage{73} - \hfill n$^{\mathrm{o}}$ 239

\S\ 4. \uline{Prolongement de sections}. \uline{Faisceaux fins}

1. \uline{Sections d'un faisceau déduit d'un préfaisceau}.

Soient $(\mathcal F_U;\rho_{VU})$ un préfaisceau sur un espace $X$, $\mathcal F$ le faisceau qui s'en déduit. On a vu au \S\ 1, n$^{\mathrm{o}}$ 7 que, moyennant les conditions $(\mathrm{PF}_{\mathrm I}(U))$ et $(\mathrm{PF}_{\mathrm{II}}(U))$, supposées vérifiées pour tout $U\in\Phi$ \uline{toute} section de $\mathcal F$ au-dessus d'un ensemble $U\in\Phi$ est de la forme $x\longrightarrow\rho_{x,U}(f_U)$, où $f_U\in\mathcal F_U$ (\S\ 1, prop.14). Nous allons voir que dans certains cas la condition $(\mathrm{PF}_{\mathrm{II}}(U))$ à elle seule est suffisante pour qu'on puisse aboutir à cette conclusion :

PROPOSITION 1. - \uline{Soient} $X$ \uline{un espace topologique}, $(\mathcal F_U;\rho_{VU})$ \uline{un préfaisceau sur} $X$ \uline{relatif à l'ensemble} $\Phi$ \uline{de toutes les parties ouvertes de} $X$, $\mathcal F$ \uline{le faisceau déduit de ce préfaisceau}. \uline{Si} $U\in\Phi$ \uline{est un sous-espace paracompact de} $X$, \uline{et si la condition} $(\mathrm{PF}_{\mathrm{II}}(U))$ \uline{est satisfaite}, \uline{toute section de} $\mathcal F$ \uline{au-dessus de} $U$ \uline{est de la forme} $x\longrightarrow\rho_{x,U}(f_U)$.

Nous démontrerons d'abord le lemme suivant :

LEMME 1. - \uline{Soient} $Y$ \uline{un espace normal}, $(U_\alpha)_{\alpha\in A}$ \uline{un recouvrement ouvert localement fini de} $Y$. \uline{Pour tout couple d'indices} $\alpha,\beta$ \uline{et tout} $y\in U_\alpha\cap U_\beta$, \uline{supposons donné un voisinage ouvert} $V_{\alpha\beta}(y)$ \uline{de} $y$ \uline{contenu dans} $U_\alpha\cap U_\beta$. \uline{Alors}, \uline{il existe une application} $y\longrightarrow V(y)$ \uline{de} $Y$ \uline{dans l'ensemble des parties ouvertes de} $Y$, \uline{telle que pour tout} $y\in Y$, \uline{on ait} $y\in V(y)$, \uline{et que les conditions suivantes soient remplies} :

1) \uline{si} $y\in U_\alpha\cap U_\beta$, $V(y)\subset V_{\alpha\beta}(y)$ ;

2) \uline{si} $y',y''$ \uline{sont tels que} $V(y')\cap V(y'')\ne\phi$, \uline{il existe un} $\alpha$ \uline{tel que} $V(y')\cup V(y'')\subset U_\alpha$.

En effet, $(U_\alpha)$ étant un recouvrement localement fini et $Y$ étant
